% Test spécifique pour vérifier l'absence de contamination QCM → question
\documentclass{nfdevoirs}

\title{Test contamination QCM-Question}

\begin{document}

\begin{devoir}{
    type=DS,
    classe={Test},
    date={Test},
    calculatrice=true,
    auteur={Test},
    correction=only
}

\begin{partie}{Test contamination}

    \begin{exercice}{Séquence problématique}

        % QCM avec correction
        \begin{qcm}{points=1, style=num}
            QCM : Combien font $3+3$ ?
            \begin{choix}
                \proposition{5}
                \proposition*{6}
                \proposition{7}
            \end{choix}
        \end{qcm}

        \begin{correction}
            $3+3 = 6$, c'est de l'addition simple.
        \end{correction}

        % Question normale avec correction
        \begin{question}{points=1}
            Question : Combien font $4+4$ ?
        \end{question}

        \begin{correction}
            $4+4 = 8$
        \end{correction}

        % Autre QCM pour vérifier reset
        \begin{qcm}{points=1, style=num}
            Autre QCM : Combien font $5+5$ ?
            \begin{choix}
                \proposition{9}
                \proposition{11}
                \proposition*{10}
            \end{choix}
        \end{qcm}

        \begin{correction}
            $5+5 = 10$, toujours de l'addition.
        \end{correction}

    \end{exercice}

\end{partie}

\end{devoir}

\end{document}